\selectlanguage{spanish}

\renewcommand{\articuloidiomaxml}{es}
\renewcommand{\articulotitulo}{Las deudas pendientes de la digitalización en América Latina}
\renewcommand{\articuloautores}{Lucas Samyn}
\renewcommand{\articulodoi}{10.18356/9789210029582c005}
\renewcommand{\articulotipo}{EDITORIAL}
\renewcommand{\articulorevista}{Blensera}
\renewcommand{\articuloissne}{2697-3650}
\renewcommand{\articulovolumen}{01}
\renewcommand{\articulonumero}{01}
\renewcommand{\articulofecha}{ 2026}

\setcounter{page}{1}
\thispagestyle{firstpage}

% --- TÍTULO A ANCHO COMPLETO ---
\noindent\begin{minipage}[t]{\dimexpr\textwidth+\marginparsep+\marginparwidth\relax}%
{\sffamily\bfseries\Large\articulotitulo\par}%
\end{minipage}\par
\vspace{3mm}%
\renewcommand{\articulotitulotrans}{The pending debts of digitalization in Latin America}
\noindent\begin{minipage}[t]{\dimexpr\textwidth+\marginparsep+\marginparwidth\relax}%
{\sffamily\itshape\normalsize\articulotitulotrans\par}%
\end{minipage}\par
\vspace{4mm}%

% --- BLOQUE LATERAL DE METADATOS ---
\begin{textblock*}{68mm}(\gbbloquex,92mm)
{\setlength{\leftskip}{0pt}\setlength{\parindent}{0pt}%
\noindent\begin{minipage}[t]{68mm}
\sffamily\small\setlength{\parskip}{1pt}
{\bfseries\color{azulrevista}Lucas Samyn}\par
\vspace{6pt}
{\bfseries\color{azulrevista!70}Derechos de autor\'ia:}\par
\textcopyright\ 2026 Samyn, L. Publicado por Blensera. Se permite el uso y distribución sin fines comerciales, siempre que se cite la obra original y no se realicen obras derivadas. (\href{https://creativecommons.org/licenses/by-nc-nd/4.0/}{https://creativecommons.org/licenses/by-nc-nd/4.0/})\par\vspace{6pt}
{\bfseries\color{azulrevista!70}Publicación:} 2026, vol.~01, n\'um.~01\par\vspace{6pt}
{\bfseries\color{azulrevista!70}DOI:} {\scriptsize\texttt{\href{https://doi.org/10.18356/9789210029582c005}{10.18356/9789210029582c005}}}\par\vspace{6pt}
{\bfseries\color{azulrevista!70}Licencia:} CC BY-NC-ND 4.0\par
{\scriptsize\texttt{\href{https://creativecommons.org/licenses/by-nc-nd/4.0/}{https://creativecommons.org/licenses/by-nc-nd/4.0/}}}\par
\vspace{6pt}
\end{minipage}%
}% FIN GRUPO leftskip
\end{textblock*}


\begin{tcolorbox}[colback=brown!8!white,colframe=brown!8!white,boxrule=0pt,arc=2pt,left=4pt,right=4pt,top=3pt,bottom=3pt,before skip=4pt,after skip=4pt]
\begin{otherlanguage}{spanish}
\sffamily\small \textbf{Resumen:} Las deudas pendientes de la digitalización en América Latina

La pandemia de COVID-19 aceleró la adopción de tecnologías en América Latina, forzando la virtualización de la educación, los trámites públicos y la salud. Sin embargo, este avance no borró las desigualdades, sino que evidenció y profundizó las fracturas sociales preexistentes en una región con una infraestructura digital muy desigual. Mientras las grandes ciudades contaban con banda ancha estable, las zonas rurales y periféricas carecían de conectividad básica, dispositivos adecuados y competencias digitales.

Tres años después de la crisis sanitaria, el panorama es ambivalente. Si bien aumentó la penetración de Internet, se consolidó la **"brecha digital de segundo nivel"** (la desigualdad en la calidad de uso). Los sectores vulnerables acceden a la red principalmente mediante teléfonos móviles con planes de datos limitados. Esto restringe su capacidad para estudiar, trabajar o ejercer su ciudadanía, actividades que requieren conexiones estables y pantallas grandes. Según datos de la CEPAL (2022), para el 20\% de la población con menores ingresos (primer quintil), el costo de la banda ancha representa en promedio el 14\% de sus ingresos mensuales, lo que imposibilita una conectividad universal real.

El desafío, por lo tanto, no es puramente tecnológico. Proveer infraestructura no basta: la digitalización sin políticas públicas integrales (que incluyan formación, diseño institucional y contenidos pertinentes) corre el riesgo de amplificar las brechas históricas de ingreso, género, etnia, territorio y edad. La transformación digital no puede entenderse aislada de las condiciones materiales de la región, y ninguna innovación tecnológica resolverá por sí sola estos problemas estructurales si no se acompaña de un debate y un diseño político profundo.
\par\smallskip\noindent\textbf{Palabras clave:} Brecha digital América latina Digitalización Políticas públicas Desigualdad social Conectividad CEPAL Inclusión digital
\end{otherlanguage}
\end{tcolorbox}

\begin{tcolorbox}[colback=brown!8!white,colframe=brown!8!white,boxrule=0pt,arc=2pt,left=4pt,right=4pt,top=3pt,bottom=3pt,before skip=4pt,after skip=4pt]
\begin{otherlanguage}{english}
\sffamily\small \textbf{Abstract:} The Unresolved Debts of Digitalization in Latin America

The COVID-19 pandemic accelerated technology adoption in Latin America, forcing the virtualization of education, public procedures, and healthcare. However, this progress did not erase inequalities; instead, it exposed and deepened pre-existing social fractures in a region with highly unequal digital infrastructure. While major cities enjoyed stable broadband, rural and peripheral areas lacked basic connectivity, adequate devices, and digital literacy.

Three years after the health crisis, the landscape remains ambivalent. Although internet penetration increased, a **"second-level digital divide"** (inequality in the quality of use) became firmly established. Low-income sectors access the network primarily via mobile phones with limited data plans. This restriction hampers their ability to study, work, or engage in civic life—activities that require stable connections and larger screens. According to CEPAL data (2022), the cost of broadband for the poorest 20\% of the population (the first quintile) represents an average of 14\% of their monthly income, making true universal connectivity impossible.

Therefore, the challenge is not purely technological. Providing infrastructure is not enough; digitalization without comprehensive public policies (including training, institutional design, and relevant content) risks amplifying historic gaps in income, gender, ethnicity, territory, and age. Digital transformation cannot be understood in isolation from the region's material conditions, and no technological innovation will solve these structural problems on its own without deep political debate and design.
\par\smallskip\noindent\textbf{Keywords:} Digital divide Latin America Digitalization Public policy Social inequality Connectivity ECLAC (Economic Commission for Latin America and the Caribbean) Digital inclusion
\end{otherlanguage}
\end{tcolorbox}

\bigskip
\noindent\rule{\linewidth}{0.4pt}
\bigskip

La pandemia de COVID-19 funcionó como un acelerador involuntario de procesos que llevaban años en gestación. En cuestión de semanas, sistemas educativos enteros migraron a plataformas virtuales, trámites administrativos que exigían presencialidad se trasladaron a portales web y la telemedicina dejó de ser una promesa para convertirse en una necesidad imperiosa. Sin embargo, esa aceleración no hizo más que volver más visibles las fracturas que ya existían en el tejido social de la región.

América Latina llegó a la pandemia con una infraestructura digital profundamente desigual. Mientras las capitales y las grandes áreas metropolitanas disponían de conexiones de banda ancha razonablemente estables, vastas zonas rurales e incluso periferias urbanas carecían de conectividad básica. La brecha no era solo de acceso físico: se manifestaba también en las competencias digitales de la población, en la disponibilidad de dispositivos adecuados y en la capacidad institucional para sostener servicios en línea de manera continua y segura.

Tres años después del momento más agudo de la crisis sanitaria, el balance resulta ambivalente. Por un lado, se registró una expansión significativa de la conectividad: varios países de la región superaron tasas de penetración de Internet que parecían lejanas a comienzos de 2020. Por otro, las desigualdades de uso —lo que la literatura especializada denomina “brecha digital de segundo nivel”— no solo persistieron sino que, en algunos casos, se profundizaron. Los sectores de menores ingresos accedieron mayoritariamente a través de teléfonos móviles con planes de datos limitados, lo que restringía severamente sus posibilidades de participar en actividades educativas, laborales o de ejercicio ciudadano que demandan conexiones estables y pantallas de mayor tamaño.\footnote{Según estimaciones de la CEPAL, en 2022 el costo del servicio de
banda ancha para la población del primer quintil de ingresos
representaba en promedio el 14\% de su ingreso mensual en la región, una
proporción que vuelve inviable cualquier pretensión de universalidad
efectiva.}

El problema, conviene insistir, no es exclusivamente tecnológico. La digitalización que no va acompañada de políticas públicas integrales —que atiendan simultáneamente la conectividad, la formación, el diseño institucional y la producción de contenidos pertinentes— corre el riesgo de reproducir y amplificar las desigualdades preexistentes. En una región donde la matriz de la desigualdad opera simultáneamente por ingreso, género, etnia, territorio y edad, suponer que la mera provisión de infraestructura resuelve el acceso es una simplificación que ningún responsable de política pública debería permitirse.

Este número de la revista reúne contribuciones que abordan distintas aristas de esta problemática. Los trabajos que aquí se presentan comparten una premisa: la transformación digital de nuestras sociedades no puede analizarse al margen de las condiciones materiales y simbólicas en las que se despliega. Esperamos que su lectura contribuya a enriquecer un debate que, lejos de cerrarse, se vuelve más urgente con cada innovación tecnológica que promete —una vez más— resolver por sí sola las deudas históricas de la región.

